% Gerado por ferramentas/gerar.py; não edite à mão.
\section*{Exercícios}\markright{Exercícios}
\addcontentsline{toc}{section}{Exercícios}
Os pontos antes de cada número indicam a dificuldade: \nivel{1} fácil, \nivel{2} médio, \nivel{3} difícil e \nivel{4} desafio. Os exercícios marcados com \estrela{} têm solução completa, com animação do gráfico, no repositório do livro: \url{https://github.com/AleffAzevedo/calculo-1-pt/blob/main/exercicios/cap3/README.md}. As respostas estão no fim do capítulo.

\begin{itemize}[leftmargin=5.2em,labelsep=0.6em,itemsep=4pt]
\item[\exnum{1}{3.1}] Use a definição de derivada para calcular $f'(x)$ para $f(x)=x^2-3x$.
\item[\exnum{1}{3.2}] Use a definição para calcular $f'(x)$ para $f(x)=\dfrac1{x+1}$.
\item[\exnum{1}{3.3}] Use a definição para calcular $f'(x)$ para $f(x)=\sqrt{2x+1}$. (Dica: conjugado.)
\item[\exnum{1}{3.4}] Encontre a reta tangente a $y=x^3$ em $x=2$.
\item[\exnum{1}{3.5}] Derive $y=4x^5-\dfrac3{x^2}+2\sqrt x-7$.
\item[\exnum{1}{3.6}] Derive $y=(x^2+1)(x^3-2x)$ de duas formas: pela regra do produto e expandindo antes.
\item[\exnum{1}{3.7}] Derive $y=\dfrac{2x-1}{x+3}$.
\item[\exnum{1}{3.8}] Derive $y=\dfrac{x^2}{x^2+1}$.
\item[\exnum{1}{3.9}] Derive $y=x^2\sen x$.
\item[\exnum{1}{3.10}] Derive $y=\dfrac{\tg x}{x}$.
\item[\exnum{1}{3.11}] Derive $y=(3x^2-5)^4$.
\item[\exnum{1}{3.12}] Derive $y=\sen(4x+1)$ e $y=\cos^2x$.
\item[\exnum{1}{3.13}] Derive $y=\sqrt{1+x^2}$.
\item[\exnum{1}{3.14}] Derive $y=e^{-3x}\cos x$.
\item[\exnum{1}{3.15}] Derive $y=\ln(x^2+4)$ e $y=\ln(\sen x)$.
\item[\exnum{1}{3.16}] Derive $y=x^3\,2^x$.
\item[\exnum{1}{3.17}] Derive $y=\arcsen(2x)$ e $y=\arctg(3x)$.
\item[\exnum{1}{3.18}] Calcule $y''$ para $y=x^4-2x^3+x$ e para $y=\sen(2x)$.
\item[\exnum{1}{3.19}] Para $s(t)=t^3-6t^2+9t$ (m), calcule a velocidade e a aceleração em $t=2$ s.
\item[\exnum{1}{3.20}] Use derivação implícita para encontrar $y'$ em $x^2+xy+y^2=7$.
\item[\exnum{2}{3.21}] \estrela\ \textbf{Retas tangente e normal.} Encontre as retas tangente e normal à parábola $y=x^2-4x+5$ no ponto de abscissa
$x=3$ e esboce as três curvas.
\item[\exnum{2}{3.22}] \estrela\ \textbf{Construindo o gráfico de $f'$.} Para $f(x)=x^3-3x^2$, determine onde $f'(x)>0$, $f'(x)<0$ e $f'(x)=0$, e esboce
o gráfico de $f'$ alinhado com o de $f$.
\item[\exnum{2}{3.23}] \estrela\ \textbf{Função por partes derivável.} Encontre $a$ e $b$ para que $f(x)=\begin{cases}x^2,&x\le1\\ax+b,&x>1\end{cases}$
seja derivável em todo $\mathbb R$.
\item[\exnum{2}{3.24}] Derive $y=(x^2+1)^3(2x-1)^2$ e fatore o resultado.
\item[\exnum{2}{3.25}] Derive $y=\sen^2(3x)$ e simplifique.
\item[\exnum{2}{3.26}] Derive $y=\cos\sqrt x$.
\item[\exnum{2}{3.27}] Derive $y=x\,e^{-x^2}$ e encontre onde $y'=0$.
\item[\exnum{2}{3.28}] Derive $y=\ln(\cos x)$ e $y=\ln|\sec x+\tg x|$.
\item[\exnum{2}{3.29}] Derive $y=\arctg\sqrt x$.
\item[\exnum{2}{3.30}] \estrela\ \textbf{O quociente de diferenças do seno.} Para $h=1$, $0{,}5$ e $0{,}1$, esboce $q_h(x)=\dfrac{\sen(x+h)-\sen x}{h}$ junto com
$\cos x$. O que acontece quando $h\to0$? Prove o resultado.
\item[\exnum{2}{3.31}] \estrela\ \textbf{Regra da cadeia e tangente.} Encontre a reta tangente a $y=\sen(x^2)$ em $x=\sqrt{\pi}$.
\item[\exnum{2}{3.32}] Calcule $y''$ para $y=xe^x$ e deduza uma fórmula para $y^{(n)}$.
\item[\exnum{2}{3.33}] Calcule $\dfrac{d^{50}}{dx^{50}}\cos x$ e $\dfrac{d^{99}}{dx^{99}}\sen x$.
\item[\exnum{2}{3.34}] Encontre os pontos em que $y=\dfrac{x}{x^2+1}$ tem tangente horizontal.
\item[\exnum{2}{3.35}] Use $f(a+h)\approx f(a)+f'(a)h$ para estimar $\sqrt{4{,}1}$ e $\sqrt[3]{8{,}3}$.
\item[\exnum{2}{3.36}] Sabendo que $g(1)=2$, $g'(1)=3$, $f(2)=5$ e $f'(2)=-1$, calcule $h'(1)$ para $h=f\circ g$ e para $k=g\cdot(f\circ g)$.
\item[\exnum{2}{3.37}] \estrela\ \textbf{Derivada de uma inversa sem a fórmula.} Seja $f(x)=x+e^x$. Mostre que $f$ é inversível e calcule $(f^{-1})'(1)$ e a reta
tangente ao gráfico de $f^{-1}$ no ponto $(1,0)$.
\item[\exnum{2}{3.38}] Use as propriedades do logaritmo antes de derivar: $y=\ln\dfrac{x^2+1}{x-1}$.
\item[\exnum{2}{3.39}] \estrela\ \textbf{Custo marginal.} O custo (em mil reais) de produzir $x$ unidades é $C(x)=0{,}01x^3-0{,}6x^2+15x+200$.
Calcule o custo marginal em $x=20$ e compare com o custo real de produzir a 21ª unidade.
\item[\exnum{2}{3.40}] Derive $y=3^{\sen x}$ e $y=\log_2(x^2+1)$.
\item[\exnum{3}{3.41}] \estrela\ \textbf{Partícula em movimento.} Uma partícula tem posição $s(t)=t^3-9t^2+24t$ (metros, $t\ge0$ em segundos).
(a) Quando ela está em repouso? (b) Quando se move para a frente e para trás?
(c) Qual é a distância total percorrida em $[0,5]$? (d) Quando está acelerando?
\item[\exnum{3}{3.42}] \estrela\ \textbf{Tangentes horizontais numa elipse.} Encontre os pontos da curva $x^2-xy+y^2=3$ onde a reta tangente é horizontal e onde é
vertical.
\item[\exnum{3}{3.43}] \estrela\ \textbf{O mínimo de $x^x$.} Use derivação logarítmica para encontrar $\dfrac{d}{dx}x^x$ ($x>0$) e determine o ponto onde
o gráfico de $y=x^x$ tem tangente horizontal.
\item[\exnum{3}{3.44}] \estrela\ \textbf{Tangentes por um ponto externo.} Encontre todas as retas tangentes à parábola $y=x^2$ que passam pelo ponto $(1,-3)$,
que não está na parábola.
\item[\exnum{3}{3.45}] Prove, pela definição, que $\dfrac{d}{dx}\cos x=-\sen x$.
\item[\exnum{3}{3.46}] Deduza $\dfrac{d}{dx}\arctg x=\dfrac1{1+x^2}$ a partir do teorema da derivada da inversa.
\item[\exnum{3}{3.47}] Derive $y=x^{\sen x}$ ($x>0$).
\item[\exnum{3}{3.48}] Use derivação logarítmica: $y=\dfrac{(x+1)^2(x-2)^3}{(x+3)^4}$.
\item[\exnum{3}{3.49}] Para a circunferência $x^2+y^2=r^2$, mostre que $y''=-\dfrac{r^2}{y^3}$.
\item[\exnum{3}{3.50}] Encontre a reta tangente à curva $x^2y+xy^2=6$ no ponto $(1,2)$.
\item[\exnum{3}{3.51}] Onde $f(x)=|x^2-4|$ é derivável? Calcule as derivadas laterais nos pontos problemáticos.
\item[\exnum{3}{3.52}] Mostre que $(\senh x)'=\cosh x$, $(\cosh x)'=\senh x$ e $(\tgh x)'=\operatorname{sech}^2x$.
\item[\exnum{3}{3.53}] Seja $f(x)=x^2\sen\frac1x$ para $x\neq0$ e $f(0)=0$. Calcule $f'(0)$ pela definição e mostre que $f'$ não é contínua em $0$.
\item[\exnum{3}{3.54}] A reta normal a $y=e^x$ em $x=0$ corta o eixo $x$ em que ponto?
\item[\exnum{3}{3.55}] Para que valor de $a$ a reta $y=ax$ é tangente à curva $y=e^x$? Em que ponto?
\item[\exnum{3}{3.56}] Derive $y=\arcsen x+\arccos x$ e conclua uma identidade.
\item[\exnum{3}{3.57}] Mostre que $y=e^{2x}+e^{-3x}$ satisfaz $y''+y'-6y=0$.
\item[\exnum{3}{3.58}] Encontre a reta tangente a $y=x^x$ em $x=e$.
\item[\exnum{3}{3.59}] O ângulo de um pêndulo é $\theta(t)=0{,}2\cos(2\pi t)$ rad. Qual é a velocidade angular máxima e quando ocorre?
\item[\exnum{3}{3.60}] Em que pontos a tangente a $y=x^3$ é paralela à reta $y=12x+1$? Escreva as tangentes.
\item[\exnum{4}{3.61}] Prove que $(fgh)'=f'gh+fg'h+fgh'$ e generalize para $n$ fatores.
\item[\exnum{4}{3.62}] Para que valores de $n$ a função $y=x^n$ satisfaz $x^2y''-3xy'+4y=0$?
\item[\exnum{4}{3.63}] Prove por indução que $\dfrac{d^n}{dx^n}(xe^x)=(x+n)e^x$.
\item[\exnum{4}{3.64}] Mostre que as curvas $x^2-y^2=a$ e $xy=b$ ($a,b\neq0$) se cortam em ângulo reto.
\item[\exnum{4}{3.65}] Encontre todas as retas que são tangentes simultaneamente a $y=x^2$ e a $y=-x^2+4x-4$.
\end{itemize}
\subsection*{Respostas}\markright{Respostas}
{\small\begin{multicols}{2}\raggedright
\begin{itemize}[leftmargin=2.7em,labelsep=0.5em,itemsep=1pt]
\item[\textbf{3.1}] $2x-3$.
\item[\textbf{3.2}] $-\dfrac1{(x+1)^2}$.
\item[\textbf{3.3}] $\dfrac1{\sqrt{2x+1}}$.
\item[\textbf{3.4}] $y=12x-16$.
\item[\textbf{3.5}] $20x^4+\dfrac6{x^3}+\dfrac1{\sqrt x}$.
\item[\textbf{3.6}] $5x^4-3x^2-2$.
\item[\textbf{3.7}] $\dfrac7{(x+3)^2}$.
\item[\textbf{3.8}] $\dfrac{2x}{(x^2+1)^2}$.
\item[\textbf{3.9}] $2x\sen x+x^2\cos x$.
\item[\textbf{3.10}] $\dfrac{x\sec^2x-\tg x}{x^2}$.
\item[\textbf{3.11}] $24x(3x^2-5)^3$.
\item[\textbf{3.12}] $4\cos(4x+1)$; $-2\sen x\cos x=-\sen2x$.
\item[\textbf{3.13}] $\dfrac{x}{\sqrt{1+x^2}}$.
\item[\textbf{3.14}] $-e^{-3x}(3\cos x+\sen x)$.
\item[\textbf{3.15}] $\dfrac{2x}{x^2+4}$; $\cotg x$.
\item[\textbf{3.16}] $2^x\big(3x^2+x^3\ln2\big)$.
\item[\textbf{3.17}] $\dfrac2{\sqrt{1-4x^2}}$; $\dfrac3{1+9x^2}$.
\item[\textbf{3.18}] $12x^2-12x$; $-4\sen2x$.
\item[\textbf{3.19}] $v(2)=-3$ m/s; $a(2)=0$.
\item[\textbf{3.20}] $y'=-\dfrac{2x+y}{x+2y}$.
\item[\textbf{3.21}] Tangente $y=2x-4$; normal $y=-\tfrac12x+\tfrac72$.
\item[\textbf{3.22}] $f'(x)=3x(x-2)$: positiva em $(-\infty,0)\cup(2,\infty)$, negativa em $(0,2)$, nula em $x=0$ e $x=2$.
\item[\textbf{3.23}] $a=2$, $b=-1$.
\item[\textbf{3.24}] $2(x^2+1)^2(2x-1)(8x^2-3x+2)$.
\item[\textbf{3.25}] $6\sen3x\cos3x=3\sen6x$.
\item[\textbf{3.26}] $-\dfrac{\sen\sqrt x}{2\sqrt x}$.
\item[\textbf{3.27}] $e^{-x^2}(1-2x^2)$; $x=\pm\tfrac1{\sqrt2}$.
\item[\textbf{3.28}] $-\tg x$; $\sec x$.
\item[\textbf{3.29}] $\dfrac1{2\sqrt x\,(1+x)}$.
\item[\textbf{3.30}] $q_h\to\cos x$; usa $\lim\frac{\sen h}h=1$ e $\lim\frac{\cos h-1}h=0$.
\item[\textbf{3.31}] $y=-2\sqrt\pi\,(x-\sqrt\pi)$.
\item[\textbf{3.32}] $y''=(x+2)e^x$; $y^{(n)}=(x+n)e^x$.
\item[\textbf{3.33}] $-\cos x$; $-\cos x$.
\item[\textbf{3.34}] $(1,\tfrac12)$ e $(-1,-\tfrac12)$.
\item[\textbf{3.35}] $2{,}025$; $2{,}025$.
\item[\textbf{3.36}] $h'(1)=-3$; $k'(1)=3\cdot5+2\cdot(-3)=9$.
\item[\textbf{3.37}] $(f^{-1})'(1)=\tfrac12$; tangente $y=\tfrac12(x-1)$.
\item[\textbf{3.38}] $\dfrac{2x}{x^2+1}-\dfrac1{x-1}$.
\item[\textbf{3.39}] $C'(20)=3$; o custo real é $C(21)-C(20)=3{,}01$.
\item[\textbf{3.40}] $3^{\sen x}\ln3\cos x$; $\dfrac{2x}{(x^2+1)\ln2}$.
\item[\textbf{3.41}] (a) $t=2$ e $t=4$; (b) frente em $[0,2)\cup(4,\infty)$, trás em $(2,4)$; (c) 28 m; (d) em $(2,3)$ e $(4,\infty)$.
\item[\textbf{3.42}] Horizontal em $(1,2)$ e $(-1,-2)$; vertical em $(2,1)$ e $(-2,-1)$.
\item[\textbf{3.43}] $y'=x^x(\ln x+1)$; tangente horizontal em $x=1/e$, com $y=e^{-1/e}\approx0{,}692$.
\item[\textbf{3.44}] $y=6x-9$ (tangente em $x=3$) e $y=-2x-1$ (tangente em $x=-1$).
\item[\textbf{3.45}] Use $\cos(x+h)=\cos x\cos h-\sen x\sen h$ e os limites do Teorema 2.7.
\item[\textbf{3.46}] $y=\arctg x\Rightarrow y'=\dfrac1{\sec^2y}=\dfrac1{1+\tg^2y}$.
\item[\textbf{3.47}] $x^{\sen x}\Big(\cos x\ln x+\dfrac{\sen x}x\Big)$.
\item[\textbf{3.48}] $y\Big(\dfrac2{x+1}+\dfrac3{x-2}-\dfrac4{x+3}\Big)$.
\item[\textbf{3.49}] $y'=-\tfrac xy$; derive de novo e use $x^2+y^2=r^2$.
\item[\textbf{3.50}] $y'=-\tfrac85$; $y=-\tfrac85x+\tfrac{18}5$.
\item[\textbf{3.51}] Em todo $x\neq\pm2$; em $x=2$: $-4$ e $4$; em $x=-2$: $-4$ e $4$.
\item[\textbf{3.52}] Derive as definições com exponenciais; use $\cosh^2-\senh^2=1$.
\item[\textbf{3.53}] $f'(0)=0$; $f'(x)=2x\sen\tfrac1x-\cos\tfrac1x$ não tem limite em 0.
\item[\textbf{3.54}] Normal $y=-x+1$; corta em $(1,0)$.
\item[\textbf{3.55}] $a=e$, no ponto $(1,e)$.
\item[\textbf{3.56}] $y'=0$, logo $\arcsen x+\arccos x=\tfrac\pi2$.
\item[\textbf{3.57}] Substitua $y'$ e $y''$; os termos se cancelam.
\item[\textbf{3.58}] $y=e^e+2e^e(x-e)$.
\item[\textbf{3.59}] $0{,}4\pi\approx1{,}26$ rad/s, quando $\theta=0$ ($t=\tfrac14+\tfrac n2$).
\item[\textbf{3.60}] $x=\pm2$: $y=12x-16$ e $y=12x+16$.
\item[\textbf{3.61}] Aplique a regra do produto a $(fg)\cdot h$; por indução, soma de $n$ termos com uma derivada em cada.
\item[\textbf{3.62}] $n=2$ (raiz dupla de $n^2-4n+4=0$).
\item[\textbf{3.63}] Base $n=1$; passo: derive $(x+n)e^x$ e obtenha $(x+n+1)e^x$.
\item[\textbf{3.64}] Inclinações $\tfrac xy$ e $-\tfrac yx$: o produto é $-1$.
\item[\textbf{3.65}] $y=0$ e $y=4x-4$.
\end{itemize}\end{multicols}}
